\documentclass[12pt]{article}
\usepackage[utf8]{inputenc}
\usepackage[T1]{fontenc}
\usepackage[brazilian]{babel}
\usepackage{geometry}
\usepackage{amsmath}
\usepackage{listings}
\usepackage{xcolor}
\usepackage{fancyhdr}
\usepackage{titlesec}
\usepackage{enumitem}

\geometry{margin=1in}
\setlength{\headheight}{15pt}
\setlength{\parindent}{0pt}
\setlength{\parskip}{4pt}

\definecolor{primary}{HTML}{1F4E79}
\definecolor{secondary}{HTML}{0B6E99}
\definecolor{lightbg}{HTML}{F4F8FC}
\definecolor{codebg}{HTML}{F7F7F7}
\definecolor{codeframe}{HTML}{D3D3D3}

\titleformat{\section}
  {\Large\bfseries\color{primary}}
  {}
  {0pt}
  {}

\titleformat{\subsection}
  {\large\bfseries\color{secondary}}
  {}
  {0pt}
  {}

\setlist[itemize]{leftmargin=1.4em,itemsep=2pt,topsep=3pt}
\setlist[enumerate]{leftmargin=1.6em,itemsep=2pt,topsep=3pt}

\pagestyle{fancy}
\fancyhf{}
\lhead{\textbf{Programação em Python}}
% \chead{Aula 06: Laços de Repetição}
\rhead{Prof. Msc. Nator Junior Carvalho da Costa}
\cfoot{\thepage}
\renewcommand{\headrulewidth}{0.4pt}

\lstset{
    language=Python,
    basicstyle=\ttfamily\small,
    backgroundcolor=\color{codebg},
    frame=single,
    rulecolor=\color{codeframe},
    keywordstyle=\color{blue},
    commentstyle=\color{gray},
    stringstyle=\color{red},
    breaklines=true,
    showstringspaces=false,
    tabsize=4
}

\begin{document}

\begin{center}
{\LARGE \textbf{Lista de Atividades - Aula 06}}\\[4pt]
{\large Laços de Repetição (\texttt{while} e \texttt{for})}\\[10pt]
\textbf{Disciplina:} Lógica de Programação em Python\\
\textbf{Professor:} Prof. Msc. Nator Junior Carvalho da Costa\\
\textbf{Semestre:} 2026.1
\end{center}

\vspace{8pt}
\colorbox{lightbg}{%
\parbox{\dimexpr\linewidth-2\fboxsep\relax}{%
\textbf{Instruções}\\
Resolva as atividades abaixo em Python. Lembre-se:
\begin{itemize}
    \item Use \texttt{while} quando não souber quantas repetições serão feitas.
    \item Use \texttt{for} quando o intervalo de repetição for conhecido.
    \item Atualize sempre a variável de controle para evitar laço infinito.
\end{itemize}
}}

\section*{Parte 1: Atividades com \texttt{while}}

\subsection*{Atividade 1: Contagem Crescente}
Leia um número inteiro \texttt{n} e exiba os números de \texttt{1} até \texttt{n}, um por linha.

\subsection*{Atividade 2: Contagem Regressiva}
Leia um número inteiro \texttt{n} e faça uma contagem de \texttt{n} até \texttt{1}. Ao final, exiba \texttt{Fim!}.

\subsection*{Atividade 3: Soma com Sentinela}
Leia números inteiros e acumule a soma até que o usuário digite \texttt{0}. Mostre a soma final.

\subsection*{Atividade 4: Senha}
Defina uma senha correta (\texttt{"python123"}) e continue pedindo a senha até o usuário acertar. Ao acertar, exiba \texttt{Acesso liberado}.

\subsection*{Atividade 5: Média de Notas}
Leia notas enquanto o usuário digitar valores maiores ou iguais a \texttt{0}. Quando ele digitar \texttt{-1}, encerre e mostre a média das notas informadas.

\section*{Parte 2: Atividades com \texttt{for}}

\subsection*{Atividade 6: Números Pares}
Use \texttt{for} para mostrar apenas os números pares de \texttt{1} até \texttt{50}.

\subsection*{Atividade 7: Tabuada}
Leia um número inteiro e mostre a tabuada desse número de \texttt{1} a \texttt{10}.

\subsection*{Atividade 8: Soma de 1 até n}
Leia um número inteiro \texttt{n} e calcule a soma de \texttt{1 + 2 + ... + n} usando \texttt{for}.

\subsection*{Atividade 9: Contador de Vogais}
Leia uma palavra e conte quantas vogais (\texttt{a, e, i, o, u}) ela possui usando \texttt{for}.

\section*{Parte 3: Escolha do Laço}

\subsection*{Atividade 10: Menu Simples}
Crie um menu com as opções:
\begin{enumerate}
    \item Somar dois números
    \item Mostrar os números de 1 a 5
    \item Sair
\end{enumerate}
Repita o menu até o usuário escolher a opção \texttt{3}.

\subsection*{Atividade 11: Jogo de Adivinhação}
Defina um número secreto (por exemplo, \texttt{7}) e peça tentativas até o usuário acertar. Ao final, exiba a quantidade de tentativas.

\section*{Exemplo Base}
\begin{lstlisting}
n = 5
while n > 0:
    print(n)
    n -= 1

print("Fim!")
\end{lstlisting}

\section*{Desafio Bônus}
\subsection*{Atividade 12: Triângulo de Asteriscos}
Leia um número inteiro \texttt{n} e exiba um triângulo de asteriscos com \texttt{n} linhas.

\begin{lstlisting}
*
**
***
...
\end{lstlisting}

\end{document}
